\documentclass{nfr}

% This LaTeX template is based on the official Word template found here: https://www.forskningsradet.no/siteassets/utlysninger/vedlegg-utlysninger/template-for-project-description-researcher-project-for-experienced-scientists-and-early-career-scientists.docx
% See the README.md file for more information

% Here's an example bibliography using the numeric (Nature) style and two
% columns for the reference list.

% \usepackage{multicol}
% \renewcommand{\bibpreamble}{\begin{multicols*}{2}}
% \renewcommand{\bibpostamble}{\end{multicols*}}

% \usepackage{biblatex}
% \addbibresource{references1.bib}

% \citestyle{nature}
% \citestyle{}
\usepackage{comment}

% \bibliographystyle{bst/veryshortnum}
\bibliographystyle{apalike}

% \title{Research Statement}
\title{Application for the Postdoctoral Research Position at Max Planck Institute for Innovation and Competition}
\author{Zhongli Wang}

\begin{document}

% Project description template
% Researcher Project for Experienced Scientists/Early Career Scientists

% This template consists of
% 
% \textbf{A:} Overview of the chapters and sections to be used as the
% structure of the project description
% 
% \textbf{B:} General guidance
% 
% \textbf{C:} Guidance on the content in each chapter and section
% 
% \textbf{A: Chapters and sections to be used}

% I used pandoc to translate the word template to tex. The formatting of the
% appendices are, therefore, not great, but the contents are there.

\maketitle

% Understanding firm competition in modern markets remains a fundamental challenge in economics. 
% While classical models such as Bertrand, Cournot, and Hotelling provide valuable benchmarks, real markets are characterized by firm asymmetries, network effects, and heterogeneous consumers, which complicate predictions and strategic decision-making \citep{Hotelling1929}.
% My research tackles this challenge by combining computational and statistical modeling to study strategic behavior and market dynamics in complex economic environments. Using tools from agent-based modeling, industrial organization, and econometrics, I analyze how firms and consumers interact in digital and networked markets. The broader goal is to bridge theory, simulation, and data to generate solutions that improve market understanding and support effective policy and managerial decisions.

% \section*{Research Statement}
\noindent Dear Members of the Recruitment Committee,\\[0.2cm]

I am writing to apply for the Postdoctoral Research Fellow position at the Max Planck Institute for Innovation and Competition. 
I am currently completing my Ph.D.\ in Economics at Bielefeld University and Università Cattolica del Sacro Cuore within the Marie Skłodowska-Curie EPOC program. 
My research lies at the intersection of innovation, digital transformation, and behavioral economic modeling. 
Using tools from machine learning, agent-based modeling, and other computational methods, I study how individuals and firms make strategic decisions in technology-intensive environments and how digital infrastructures, information frictions, and regulatory conditions shape innovation dynamics. 
These themes align closely with the research agenda on the determinants, mechanisms, and societal implications of innovation processes, as well as with its interest in the behavioral foundations of inventive and entrepreneurial activity.\\

I began my PhD in 2021, at a time when digitalization, data-driven technologies, and artificial intelligence were reshaping how firms innovate and how new technologies diffuse through markets. 
This context motivated the core projects of my dissertation, which examine how individual and organizational incentives, technological constraints, and institutional environments jointly shape innovation dynamics and strategic behavior. 
Using machine learning, agent-based modeling, and simulation-based econometric methods, I study how heterogeneous agents interact within networked environments, how digital infrastructures influence the development and adoption of new technologies, and how regulatory interventions affect innovation incentives and market outcomes. 
These projects required designing behavioral models of inventive activity, estimating interaction patterns, and simulating large-scale innovation and diffusion processes. 
All three chapters of my dissertation are single-authored, underscoring my ability to independently conceptualize and execute complex computational research projects from start to finish.\\

A central theme of my research is understanding how individual and organizational behavior evolves within networked and technologically dynamic environments, and how these interaction patterns shape large-scale innovation outcomes. 
Although my dissertation focuses on digital markets, the underlying mechanisms I analyze—heterogeneous incentives, learning dynamics, diffusion processes, and the emergence of collective patterns—are fundamental to understanding innovation and the adoption of new technologies. 
To study these dynamics, I develop agent-based models in which firms and individuals update beliefs, adjust strategic choices, and respond to technological and institutional conditions based on their network position and local information. 
I complement these models with simulation-based estimation and machine-learning techniques to quantify behavioral parameters and evaluate how alternative regulatory or informational environments influence innovation trajectories and market evolution. 
This combination of behavioral modeling, network analysis, and computational simulation provides a strong methodological foundation for investigating the determinants of innovative activity, the diffusion of new technologies, and the behavioral underpinnings of creativity and strategic adaptation.\\

In addition to my dissertation work, I have coauthored two interdisciplinary papers that demonstrate my ability to integrate modern machine-learning techniques with rigorous empirical and computational analysis \citep{Zhao2025, Dawid2025}. 
The first, \textit{“Polytope Fraud Theory”}, applies machine-learning and statistical methods to detect patterns and anomalies in large-scale financial data, illustrating my experience with complex data processing, feature extraction, and the identification of latent behavioral structures—capabilities that are increasingly valuable for analyzing inventive activity, entrepreneurial dynamics, and innovation-related datasets. 
The second, \textit{“Agentic Workflows for Economic Research”}, develops a framework for incorporating large-language-model agents into empirical workflows, enabling automated text analysis, document processing, and reproducible computational pipelines. 
Together, these projects highlight my ability to combine advanced ML tools with structured research design, and they reflect my broader interest in how digital technologies—including AI—transform analytical capabilities, organizational processes, and the broader innovation ecosystem.\\

I am particularly motivated to join the Max Planck Institute for Innovation and Competition because of its global reputation for rigorous, conceptually grounded, and methodologically innovative research on the determinants and consequences of innovation. 
The Institute’s focus on understanding inventive activity, the behavioral foundations of creativity and strategic decision-making, and the implications of digitalization and AI aligns closely with my own research agenda. 
I am especially drawn to MPIIC’s emphasis on independent basic research within a highly collaborative and interdisciplinary environment, where scholars from economics, management, and related fields engage deeply with both theoretical frameworks and empirical challenges. 
The opportunity to contribute to research on innovation dynamics—while interacting with leading researchers and accessing a vibrant intellectual community—offers an ideal setting to further develop my computational, behavioral, and analytical approaches to studying technological change and organizational adaptation. 
I would welcome the opportunity to bring my expertise in machine learning, agent-based modeling, and quantitative analysis to the Institute, and to advance research that deepens our understanding of innovation processes and their broader societal implications.\\

At the Max Planck Institute for Innovation and Competition, I hope to contribute my quantitative, computational, and behavioral modeling expertise to research on the determinants of innovation, the diffusion of new technologies, and the behavioral foundations of inventive and strategic activity. 
My areas of special interest include:
\begin{itemize}
    \item Modeling how heterogeneous individuals and organizations generate, adopt, and adapt to new technologies using agent-based and computational microeconomic approaches;
    \item Applying machine-learning and simulation techniques to study innovation-related datasets, digital knowledge production, and the evolution of technological domains;
    % \item Developing dynamic simulations to analyze how network structures, incentives, and institutional conditions influence innovation trajectories and the emergence of novel organizational or market patterns;
    \item Investigating how regulatory frameworks, information environments, and digital infrastructures shape inventive activity, innovation behavior, and the strategic decisions of firms and entrepreneurs.
\end{itemize}

I am especially motivated by the Max Planck Institute for Innovation and Competition’s commitment to rigorous, theory-driven, and methodologically innovative research on the processes underlying innovation and creative activity. 
% The Institute’s emphasis on independent basic research, combined with its interdisciplinary environment and strong focus on digital transformation and the behavioral foundations of inventive behavior, provides an ideal setting for advancing my computational and analytical work. 
% I would welcome the opportunity to contribute my expertise in machine learning, agent-based modeling, and quantitative analysis to ongoing research at the Institute, and to develop new projects that deepen our understanding of technological change, innovation incentives, and the evolution of digital and knowledge-based economies.
\\

Thank you so much for considering my application! I would welcome the opportunity to discuss how my computational and behavioral modeling expertise can contribute to the Institute’s research on innovation, technological change, and the behavioral foundations of inventive and strategic activity.
\vspace{5mm}

\noindent Sincerely,\\
\vspace{5mm}

\noindent Zhongli Wang

% \vspace{10mm}

% Throughout my doctoral and collaborative research, I have experienced the value of combining analytical modeling with computational innovation. Working across fields has shown me that the most pressing problems—such as market transparency, sustainable technology adoption, and responsible data governance—can only be solved through interdisciplinary collaboration.

% In addition to my market-oriented research, my coauthored publication \textit{``Polytope Fraud Theory''} and the working paper \textit{``Agentic Workflows for Economic Research''} have shaped my commitment to methodological innovation. The former integrates economic reasoning with machine learning to enhance fraud detection and investor protection; the latter connects computational economics with artificial intelligence by designing LLM-based research workflows. These experiences have strengthened my belief that economists can contribute not only to theoretical insights but also to the design of intelligent, transparent, and reproducible systems.

% I am therefore particularly excited about joining the FINATRAX Research Group, where business research meets information systems engineering. The group’s mission—to explore digital identities, distributed ledgers, machine learning, and IoT as enablers of sustainable technologies—closely aligns with my research vision. I am enthusiastic about participating in the design, implementation, and theoretical elaboration of high-quality, interdisciplinary solutions that create tangible value for both organizations and society. This position represents a natural continuation of my efforts to bridge economics with emerging digital technologies.

% \section*{Funding and Collaboration Strategy}

% My strategy for research funding and partnerships builds on a dual foundation of analytical rigor and applied collaboration. 
% Methodologically, I plan to combine analytical and computational approaches: analytical derivations will establish benchmark propositions, while numerical simulations and agent-based modeling will explore market dynamics and firm behavior under varying degrees of information transparency. This structure naturally supports both theoretical innovation and data-driven applications—key elements for competitive research proposals.

% For empirical validation, I will leverage access to the Refinitiv database through collaboration with the Southern University of Science and Technology (SUSTech), complemented by additional data resources available at the host institution. My doctoral research, funded by the \textit{Marie Skłodowska-Curie Actions} and the \textit{Bielefeld Young Researchers’ Fund}, has already demonstrated my ability to manage funded projects and coordinate cross-institutional collaboration. 

% Going forward, I intend to apply for European and national-level grants such as \textit{Horizon Europe}, \textit{DFG}, and \textit{VolkswagenStiftung}, focusing on projects that integrate computational economics with digital transformation research. On the industry side, I aim to build partnerships with firms in the ICT and FinTech sectors to co-develop decision-support tools and simulation frameworks for market design and policy evaluation. These collaborations would ensure practical relevance, foster technology transfer, and open opportunities for joint funding applications. With adequate computational resources and institutional support, I am confident that these efforts can yield both high-impact publications and meaningful contributions to industry and policy communities.





\begin{comment}

\section*{Statement of Research Interests}

% Here's a citation to \cite{texbook}. Here are citations to \cite{texbook,knuth:1984,lesk:1977}, and to \cite{lesk:1977,texbook,latex2e,latex:companion}.

My doctoral dissertation, ``\textit{Strategic Behavior and Economic Dynamics in ICT Markets: An Agent-based Framework}", investigates how firms in the information and communications technology (ICT) industry adapt under constraints such as data regulations, capacity limits, and consumer network effects. The work builds a bridge between state-space and agent-based models, showing how techniques from dynamic latent-variable systems can be applied to simulate and estimate complex economic environments.\footnote{All three chapters are single-authored papers.}\\

My job market paper, \textit{``Strategic Firm Capacity Selection under Network Effects: A Computational Model of Capacity Investment Dynamics in the ICT Industry''}, investigates how digital firms determine optimal service capacity under the joint influence of consumer congestion and network effects---an interaction that has received limited attention in the existing literature. While traditional industrial organization models emphasize innovation or price competition, many digital firms face a different constraint: their profitability depends critically on infrastructure investment and the ability to maintain service quality. To address this challenge, I combine game-theoretic reasoning with agent-based simulations. The game-theoretic framework captures the multi-stage nature of firms’ decisions on pricing and capacity investment, while the agent-based model represents the market dynamics that emerge when heterogeneous consumers interact through a small-world network and share information within user communities. The results show that strong network effects tend to amplify competitive advantages by concentrating demand, whereas congestion generates counteracting pressures that reduce pricing flexibility and erode profits. Together, these forces explain why digital firms must balance capacity expansion and user experience to sustain long-term profitability. This study is among the first to systematically analyze the joint role of network and congestion effects in the ICT industry, offering new insights into pricing, investment, and regulatory strategies for digital service markets.\\

Building on the similar market context, my previous paper, \textit{``Heterogeneous Consumers, Firm Competition, and the Data Value Chain: An Agent-Based Approach''}, investigates how differences in data-governance intensity influence competition in software markets. The challenge arises because digital firms depend on user data to enhance product quality and pricing decisions, while local governments impose varying levels of data regulation to protect privacy and maintain market fairness. 
To analyze this tension, I develop a hybrid game-theoretic and agent-based model in which governments determine the stringency of data governance, and firms subsequently compete on investment and pricing decisions. 
Consumers are heterogeneous in their sensitivity to product quality and their attitudes toward data governance, which shape their purchasing behavior. 
The results show that stricter data governance limits firms’ ability to utilize consumer data effectively and reduces competitive differentiation, which leads to higher prices and lower consumer welfare. In contrast, moderate levels of regulation can maintain competition and promote more balanced market outcomes. The analysis suggests that regional governments would benefit from greater coordination in setting data policies and from incorporating consumer preferences into policy design, thereby supporting both market efficiency and innovation. Overall, these findings offer a structured perspective on how policymakers can balance privacy protection with incentives for innovation in digital markets.\\

Building on the insights from the first two studies, my third paper ``\textit{Parameter Estimation Methods in an Agent-based Computational Model: Insights on ICT Industrial Dynamics and Policies}" focuses on empirical validation. 
Using financial and operational data from two leading U.S. telecommunication firms, AT\&T and T-Mobile, I estimate the model parameters with the Simulated Method of Moments and neural network–based surrogate modeling (both within-sample and out-of-sample check). 
The main challenge lies in linking complex agent-based models to real-world data, as repeated simulations are computationally intensive and traditional estimation methods often fall short. 
To address this, I employed the Paderborn High-Performance Computing (HPC) clusters and the Python package \texttt{ray} to parallelize the simulations and collect 40{,}000 samples for parameter search. 
By comparing these two approaches, the study provides practical evidence on how to calibrate computational models efficiently while maintaining empirical accuracy. This work contributes to the broader methodological discussion on integrating simulation-based models with data-driven analysis to better understand market behavior and policy outcomes.



\section*{Publications and Emerging Projects}

In addition to my market-oriented research, I have one peer-reviewed publication and one current working paper that contribute to the broader use of computational methods in economics.\footnote{Both papers are coauthored with other researchers.} \\

The paper \textit{``Polytope Fraud Theory''} extends the classical Fraud Triangle and Fraud Diamond frameworks by introducing a multidimensional structure that captures ten abnormal financial practice alarms indicative of potential corporate fraud \citep{Zhao2025}. Using data on activist short sellers and supervised machine learning techniques, the study identifies financial patterns associated with fraudulent reporting and demonstrates how algorithmic approaches can support market transparency and investor protection. 
The paper also demonstrates the potential for automating advanced financial statement analysis, providing tools that can enhance both auditing efficiency and investment decision-making. In addition, it introduces the \textit{Unified Investor Protection Framework}, which systematically categorizes investor-protection theories from the macro- to the micro-level. This framework offers practical guidance for identifying firm-level financial risks and strengthens the connection between analytical modeling and real-world investor decision processes.\\

Another working paper, \textit{``Agentic Workflows for Economic Research: Design and Implementation''}, develops a framework that employs large-language-model (LLM) agents to support economists throughout the research process—from data collection and coding to simulation design and interpretation \citep{Dawid2025}. The framework formalizes how autonomous AI agents can collaborate within multi-stage research pipelines while ensuring reproducibility, version control, and transparent documentation. Using examples implemented in Microsoft AutoGen, we demonstrate how such workflows can accelerate tasks like research-question ideation, literature review, and model specification, thereby improving research efficiency without compromising academic rigor. Importantly, the framework positions human researchers at the core of the process: AI agents act as supportive tools that augment, rather than replace, human judgment and creativity. This line of work extends my methodological interests in AI-enhanced modeling, bridging computational economics with artificial-intelligence research and tool development.

\section*{Future Research}

\subsection*{Potential Research Direction: The Observability of Firm Heterogeneity in Price-Directed Consumer Search}

% Recent models of price-directed search \citep{obradovits2023price} highlight how consumers allocate their attention and search effort across differentiated firms when facing information frictions. However, these frameworks often assume that firm heterogeneity---in cost structures, quality, or service intensity---is either fully observable to consumers or symmetrically distributed across sellers. In many real markets, such as digital platforms, online retailing, or app ecosystems, this heterogeneity is only partially observable: consumers infer firm characteristics indirectly from prices, reputation indicators, or platform signals.

% Understanding how the observability of firm heterogeneity shapes consumer search behavior and market efficiency is crucial for both theory and policy. It determines the extent of equilibrium price dispersion, the efficiency of market matching, and firms’ incentives to disclose information. The problem is empirically feasible because modern datasets (e.g., online pricing panels, platform APIs, or laboratory experiments) provide micro-level evidence on search behavior under varying degrees of information transparency. Computational approaches, including agent-based simulations, further enable the exploration of equilibrium dynamics in complex environments.

Recent models of price-directed search \citep{obradovits2023price} examine how consumers allocate attention and search effort across differentiated firms under information frictions, yet they typically assume that firm heterogeneity---in costs, quality, or service intensity---is fully observable or symmetrically distributed. In real markets such as digital platforms or online retailing, this heterogeneity is only partially observable, as consumers infer firm characteristics indirectly from prices, ratings, or platform signals. Understanding how such observability affects consumer search and market efficiency is crucial, as it shapes equilibrium price dispersion, allocative outcomes, and firms’ incentives to disclose information. The question is empirically feasible given modern micro-level data on online search and pricing, while computational approaches, including agent-based simulations, allow systematic analysis of equilibrium dynamics under varying information transparency. \\

% This research direction contributes to Industrial Organization by linking information structure, firm heterogeneity, and price competition under bounded rationality. It combines theoretical modeling with simulation and empirical validation, making it both conceptually rich and methodologically tractable. 

The following hypotheses will be tested:

\begin{enumerate}%[label=\textbf{H\arabic*.}, leftmargin=2em]
    \item \textbf{Partial observability increases price dispersion.} When firm heterogeneity (e.g., in cost or quality) is imperfectly observable, consumers engage in less directed search, resulting in a wider equilibrium price distribution than under full observability.
    
    \item \textbf{High observability improves allocative efficiency.} Markets with more transparent signals of firm heterogeneity (such as ratings or quality certifications) achieve a closer match between consumers and preferred firms, raising consumer surplus and reducing welfare loss.
    
    \item \textbf{Strategic transparency depends on firm type.} Low-cost or high-quality firms have stronger incentives to increase the observability of their characteristics (e.g., through signaling or advertising) than high-cost or low-quality competitors, leading to asymmetric information equilibria.
\end{enumerate}

\subsection*{Methodological Plan}

In my approach, I plan to combine analytical and computational methods. Analytical derivations will be used to establish the basic propositions and benchmark outcomes of the model, while numerical techniques and computational modeling will be employed to solve for Nash equilibria and explore market dynamics under varying degrees of information transparency. For empirical validation, I have access to the Refinitiv database through an application via the Southern University of Science and Technology (SUSTech) and may also rely on additional database resources at my future host institute. My doctoral research was supported by the Marie Skłodowska-Curie Actions and the Bielefeld Nachwuchsfonds, and it benefited from extensive computational work conducted on Google Colab (self-funded) and the Paderborn High Performance Computing cluster (supported by Bielefeld University). I hope that the new institute could provide adequate support for research infrastructure, particularly computational resources, to further advance this line of study.



% \subsection*{Computational \& ABM Dynamics}

% Specifically, I plan to:

% •	Develop scalable frameworks for data-driven research of economic models that can inform competition and innovation policy.

% •	Study how AI and automation transform firm strategy, market structure, and industrial dynamics.

% My long-term objective is to contribute to the economics literature, that is both analytically rigorous and computationally adaptive—one that captures the rich, evolving behavior of firms and consumers in real markets. I remain deeply motivated to pursue research that connects theoretical innovation with practical relevance, and I look forward to collaborating in an environment that values interdisciplinary approaches and empirical credibility.
% \citep{arve2025decentralized}.


% \bibliography{NFR_Template_Statements/references}

\newpage

\appendix
\end{comment}
\bibliography{NFR_Template_Statements/references1}

\end{document}
